\section{Related work}
\label{sec:related-work}

\subsection{Jerk-limited online trajectory generation}

% CLAIM: C01,C02,C11
Jerk-limited trajectory generation has a long lineage in real-time robot
motion.  Polynomial and synchronized constructions bound higher derivatives
while moving between specified states
\citep{MacfarlaneCroft2003JerkBounded,HaschkeEtAl2008Online}.
Discrete-time filtering provides another route to trajectories with bounded
velocity, acceleration, and jerk
\citep{GerelliLoBianco2010DiscreteFilter}, while mathematical programming can
optimize a minimum-jerk criterion under constraints
\citep{CanaliEtAl2014MinimumJerk}.  Path-accurate schemes address preservation
of a known desired path under jerk limits
\citep{LangeAlbuSchaeffer2016PathAccurate}.  These objectives should not be
collapsed: terminal-state interpolation, path preservation, minimum-jerk
optimization, and following an incompletely observed moving reference are
different problems.

\Ruckig{} generates time-optimal jerk-limited trajectories to arbitrary target
states and exposes a real-time state-to-state update interface
\citep{BerscheidKroeger2021Ruckig}.  That capability is the native follower in
our ordinary state-to-state conditions.  The distinction developed here lies
upstream and at the execution boundary: a position stream does not itself
define a time-aligned target state, and the legality of the executed
control-period prefix depends on its full piecewise profile.  We make no
priority claim for either observation.

\subsection{Reference governance and predictive motion}

% CLAIM: C02,C09
Reference and command governors are add-on mechanisms that modify references
to enforce constraints for a specified plant or closed-loop model
\citep{KolmanovskyEtAl2014GovernorTutorial,GaroneEtAl2017GovernorSurvey}.
This perspective motivates the term \emph{executable-target governor}: the
layer is neither an estimator nor a replacement for the native trajectory
solver.  Our construction is substantially narrower than general governor
theory.  It uses uncoupled sampled triple-integrator command dynamics and a
one-period viability check; it does not inherit robust plant-level,
disturbance, or stability guarantees from the broader literature.

Receding-horizon optimization can combine target updates, jerk-input
kinematics, and constraints over a longer prediction window
\citep{GhazaeiArdakaniEtAl2015MPC}.  Such a formulation can trade tracking,
effort, and future feasibility jointly.  The present method instead selects a
single adjacent jerk from analytically constructed candidate intervals and
checks the resulting next state.  It is intended to make the timing and
command contract explicit, not to claim MPC-equivalent optimality.

\subsection{Position-only state estimation and timing}

% CLAIM: C05
Derivatives from sampled positions can be obtained by local-polynomial
smoothing and differentiation \citep{SavitzkyGolay1964}, fixed-gain
\(\alpha\)--\(\beta\)--\(\gamma\) trackers
\citep{Kalata1984TrackingIndex,SahoMasugi2015ABG}, or state-space filtering and
prediction \citep{Kalman1960Filtering}.  Recent work on real-time numerical
differentiation emphasizes the distinction between estimates that use data
available at the estimated time and noncausal methods that consume future
samples \citep{VermaEtAl2024Differentiation}.  Our experiments do not propose
a new estimator.  They use simple derivative sources as controlled
information conditions and retain both represented time and availability time
through the downstream pipeline.

This bookkeeping matters because filtering delay and prediction horizon
compose.  A posterior representing \(t_{k-1}\) but available at \(t_k\) must
first be propagated to the control time before a prediction \(h\) beyond
\(t_k\) has its stated meaning.  Relabeling a delayed posterior as current
state does not create information; a model-based propagation must be explicit.
Similarly, a centered stencil aligned at \(t_k\) is an offline condition if it
requires \(p_{k+1}\).

\subsection{Moving-reference interfaces and method identity}

% CLAIM: C02,C12
The official \RuckigTrackingInterface{} documentation describes an interface
that predicts ahead to reduce lag when following a sampled signal, and the API
reference identifies the corresponding \TrackigAPI{} class
\citep{RuckigTrackingTutorial,RuckigTrackigAPI}.  We cite these sources only to
establish that moving-signal following is recognized as a distinct interface
behavior.  The interface was not measured in this project; the evidence does
not show that it is necessary, superior, or available in every tested
binding.

\begin{table}[t]
  \centering
  \small
  \caption{Conceptual method taxonomy.  Target information and the actually
  executed follower are independent factors.  ``Exact prefix'' means exact
  segment data are exposed; otherwise internal jerk remains unavailable.}
  \label{tab:method-taxonomy}
  \setlength{\tabcolsep}{4pt}
  \begin{tabular}{p{0.20\linewidth}p{0.18\linewidth}p{0.20\linewidth}p{0.27\linewidth}}
    \toprule
    Method role & Reference information & Executed profile & Identity boundary \\
    \midrule
    Ordinary state-to-state OTG & terminal P/PV/PVA request & native
    piecewise prefix & unshielded native command \\
    Tracking-aware follower & sampled/predicted moving reference & interface
    dependent & not evaluated here \\
    Viability-shielded OTG & governed or native candidate & native or
    replacement & shield decision changes identity \\
    Direct governor execution & future objective plus current state & one
    exact constant-jerk segment & governor action is the command \\
    \bottomrule
  \end{tabular}
\end{table}

The synthesis in \cref{tab:method-taxonomy} is the gap this paper addresses:
the same P/PV/PVA label may feed different governors or followers, and the same
declared follower may be replaced on execution.  Consequently, a defensible
comparison needs both an information condition and an executed-method
identity.  The following formulation makes those two axes, and every relevant
clock, explicit.
